
\documentclass[aspectratio=169]{beamer}

\mode<presentation>
\usetheme{Boadilla}
\definecolor{NTHU1}{RGB}{127,16,132}
\definecolor{NTHU2}{RGB}{228,210,231}
\definecolor{blue}{RGB}{30,90,205}
\definecolor{red}{RGB}{213,94,0}
\definecolor{green}{RGB}{0,128,0}
\definecolor{charcoalgray}{RGB}{108,104,104}
\setbeamercolor{title}{fg=NTHU1}
\setbeamercolor{frametitle}{fg=NTHU1}
\setbeamercolor{block title}{bg=NTHU1, fg=white}
\setbeamercolor{block body}{bg=NTHU2}
\setbeamercolor{structure}{fg=NTHU1}
\setbeamercolor{item projected}{fg=white}
\setbeamercolor{item}{fg=NTHU1}
\setbeamercolor{subitem}{fg=NTHU1}
\setbeamercolor{section in toc}{fg=NTHU1}
\setbeamercolor{description item}{fg=NTHU1}
\setbeamercolor{caption name}{fg=NTHU1}
\setbeamercolor{button}{bg=NTHU1, fg=white}
\setbeamercolor{caption name}{fg=NTHU1}
\usepackage{graphics}
\usepackage{tikz}
\usepackage{amsmath}
\usepackage{bbm}
\usetikzlibrary{decorations.pathreplacing}
\usepackage{geometry}
\usepackage{booktabs}
\usepackage{bookmark}
\usepackage{multirow, makecell}
\usepackage{float}
\usepackage{fancyvrb}
\usepackage{caption}
\captionsetup[figure]{singlelinecheck = false, format= hang, justification=centering}
\captionsetup[subfigure]{singlelinecheck = false, format= hang, justification=raggedright}
\usepackage{subcaption}
\usepackage{adjustbox}
\usepackage{hyperref}
\usepackage{threeparttable}
\usepackage{appendixnumberbeamer} 
\usepackage[T1]{fontenc}
\usepackage[default]{lato} 
\usepackage{smartdiagram}
\smartdiagramset{
  back arrow disabled = true,
  module minimum width=1cm,
  module x sep = 3,
  uniform arrow color=true,
}
\usepackage{xcolor}
\usepackage{colortbl}
	\definecolor{light-gray}{gray}{0.99}

\newenvironment{wideitemize}{\itemize\addtolength{\itemsep}{10pt}}{\enditemize}
\newenvironment{wideenumerate}{\enumerate\addtolength{\itemsep}{10pt}}{\endenumerate}
\newenvironment{widedescription}{\description\addtolength{\itemsep}{10pt}}{\enddescription}
\hypersetup{
colorlinks=true,
linkcolor=NTHU1,
filecolor=green, 
urlcolor=blue,
}
\beamertemplatenavigationsymbolsempty
\setbeamercolor{author in head/foot}{bg=white, fg=NTHU1}
\setbeamercolor{title in head/foot}{bg=white, fg=NTHU1}
\setbeamercolor{date in head/foot}{bg=white, fg=NTHU1}
\setbeamercolor{section in head/foot}{bg=white, fg=NTHU1}
\setbeamercolor{headline}{bg=white}
\setbeamertemplate{footline}{
    \leavevmode%
    \hbox{%
        \begin{beamercolorbox}[wd=.333333\paperwidth,ht=2.25ex,dp=1ex,center]{date in head/foot}%
            \usebeamerfont{date in head/foot}%\insertdate
        \end{beamercolorbox}%
        \begin{beamercolorbox}[wd=.333333\paperwidth,ht=2.25ex,dp=1ex,center]{title in head/foot}%
            \usebeamerfont{title in head/foot}\insertshorttitle
        \end{beamercolorbox}%
        \begin{beamercolorbox}[wd=.333333\paperwidth,ht=2.25ex,dp=1ex,right]{date in head/foot}%
            \usebeamerfont{date in head/foot}\insertframenumber{}\hspace*{2em}
        \end{beamercolorbox}}%
        \vskip0pt%
    }
    %% May need to script the introduction!!!!!!
%\setbeamercolor{page number in head/foot}{fg=NTHU1}
\setbeamertemplate{section in toc}[sections numbered]
\setbeamertemplate{subsection in toc}{\leavevmode\leftskip=3em\rlap{\hskip-1.75em\inserttocsectionnumber.\inserttocsubsectionnumber}
\inserttocsubsection\par}
\setbeamerfont{subsection in toc}{size=\footnotesize}
%\setbeamertemplate{headline}{%
  %\begin{beamercolorbox}[ht=5.5ex]{section in head/foot}
    %\vskip2pt\insertnavigation{0.33\paperwidth}\vskip2pt
  %\end{beamercolorbox}%
%}
\newenvironment{transitionframe}{\setbeamercolor{background canvas}{bg=white}\setbeamertemplate{footline}{} \begin{frame}[noframenumbering]}{\end{frame}}


\makeatletter
\let\@@magyar@captionfix\relax
\makeatother


\title[]{Public Economics and Development} % Change this regularly
\subtitle[]{Week 7: Challenges to taxation in weak states: Informal economies and political economy}
\author[Lee]{Seung-hun Lee }
\institute[]{Taipei School of Economics\\ National Tsing Hua University}

\date[]{April 14th, 2026}

\begin{document}
\begin{transitionframe}
\titlepage
\end{transitionframe}


%%% Color slides for section headers: Use for colloquium version (The ones bewteen \iffals and \fi)

\section{Overview setup of the topic}
\begin{frame}
\frametitle{Development of fiscal capacity is a key feature of development} 
Expansion of government size and fiscal capacity investment in the last century 
\begin{itemize}\setlength\itemsep{3pt}a
\item Tax share of GDP: 10\% in 1910 $\to$ 30-50\% today
\item Tax capacity improved following implementation of broad-based taxes (income tax)
\item However, this explains the trend only in developed countries
\begin{itemize}\setlength\itemsep{3pt}
\item So economic growth $\to$ better fiscal capacity? Not regularly 
\item Separate investment in capacities are key
\end{itemize}
\end{itemize}\medskip
But in poor states, no noticeable increase in fiscal capacities
\begin{itemize}\setlength\itemsep{3pt}
\item Raise significantly less revenues than developed countries
\item Lag behind in the implementation of modern tax technologies
\begin{itemize}\setlength\itemsep{3pt}
\item Property registers are manually updated, if updated at all 
\item Do not utilize third-party information that ensures accurate tax information
\item Less digitalized systems, more loopholes in general
\end{itemize}
\end{itemize}
\end{frame}


\begin{frame}
\frametitle{Less capable states rely more on distortionary and narrow based taxes} 
Why do they depend on such taxation
\begin{itemize}\setlength\itemsep{3pt}
\item Broad-base taxes require registries that cover formal sector, and audits
\item Narrow-base taxes rely on easily observable bases (trade ports, mines etc)
\item Large informal economy and underdeveloped finance sector: Opens up loopholes
\item Extractive for select groups by shifting tax burden onto them
\item Bad equilibrium: Fiscal capacity$\Downarrow\to$ Public goods and economic development $\Downarrow$
\end{itemize}\medskip
What the data tells us
\begin{itemize}\setlength\itemsep{3pt}
\item High income countries rely on broad-based taxes
\item Narrow base taxes (tariffs) more observable in low-capacity countries
\end{itemize}
\end{frame}

\begin{frame}
\frametitle{High income $\to$ Income tax, Lower income $\to$ Trade tax} 
\begin{figure}
  \includegraphics[keepaspectratio, width=0.7\textwidth]{fig1.png}
\end{figure}
\end{frame}


\begin{frame}
\frametitle{Less capable states suffer from negative shocks from conflicts} 
External Conflicts and taxation (Besley, Persson 2010)
\begin{itemize}\setlength\itemsep{3pt}
\item External wars raise value of public goods (e.g. national defense)
\item Hence, motivates increases in fiscal capacity
\end{itemize}\medskip
However, many of the global conflicts are internal
\begin{itemize}\setlength\itemsep{3pt}
\item Approx 20-30\% of countries worldwide have active civil conflicts during 1960-2000s
\item Supermajority of conflicts are intrastate (Ray,Esteban 2017)
\item Depletes administrative infrastructure used to mobilize resources
\item Hence, leads to decline in tax collection and public goods provision
\end{itemize}
\end{frame}

\begin{frame}
\frametitle{External wars increase capacity; internal wars hamper it}
\begin{figure}
  \begin{subfigure}{0.49\textwidth}
  \includegraphics[keepaspectratio,width=\textwidth]{fig2_1.png}
\caption{External wars positively correlate}  
\end{subfigure}
  \begin{subfigure}{0.49\textwidth}
  \includegraphics[keepaspectratio,width=\textwidth]{fig2_2.png}
  \caption{Internal wars negatively correlate}
  \end{subfigure}
\end{figure}
\end{frame}


\begin{transitionframe}
\frametitle{Roadmap for today}
\tableofcontents
\end{transitionframe}








\section{Theories on tax leakages}
\subsection{Allingham, Sandmo (1972): Income Tax Evasion: A Theoretical Analysis}
\begin{transitionframe}
\begin{center}
  \begin{huge}
    \textcolor{NTHU1}{%
Allingham, Sandmo (1972):\\ Income Tax Evasion: A Theoretical Analysis
    }
  \end{huge}
\end{center}
\end{transitionframe}

\begin{frame}
\frametitle{Q: How and why do some people avoid taxes?}
Optimal taxation from the standpoint of the individuals
\begin{itemize}\setlength\itemsep{3pt}
\item Up until now, we assumed that individuals abide by taxes levied on them
\item In reality, that is not always the case
 \begin{itemize}\setlength\itemsep{3pt}
\item Tax avoidance: Exploit legal loopholes to minimize tax liability
\item Tax evasion: Misreport tax liability illegally
\item The distinction between the two is \textit{subtle}...
\end{itemize}
\end{itemize}\medskip
This paper provides a classical model of tax evasion
\begin{itemize}\setlength\itemsep{3pt}
\item Individuals factor in the detection probability to decide on how much income to \textit{report}
\item The true income of each individual is given
\item If not caught, they keep more income than they would have if they complied to taxes
\item If caught, they pay taxes and penalty on difference between true and reported income
\end{itemize}
\end{frame}

\begin{frame}
\frametitle{Setting up the framework: Individual decisions} 
Assumptions
\begin{itemize}\setlength\itemsep{3pt}
\item Individual's true income is $W$
\item Tax rate $t$ levied on \textit{reported} income $\bar{W}=W-E$, where $E$ is amount of evasion
\item Government conducts tax audit with probability $p$ 
 \begin{itemize}\setlength\itemsep{3pt}
\item We assume audits always catch evasion 
\item If they do not, we can use \textit{effective} probability of detection $p'<p$
\end{itemize}
\end{itemize}\medskip
Possible scenarios on income
\begin{itemize}\setlength\itemsep{3pt}
  \item If not caught: Individuals keep 
  \[
  \begin{aligned}
  Y_N &= W-t\bar{W} = W-t(W-E)\\
  &=(1-t)W+tE\\
  \end{aligned}
  \]\vspace{-10pt}
  \item If caught: Pay taxes on $E$ at rate $\theta>t$
  \[
  \begin{aligned}
  Y_A&= W-t\bar{W} -\theta (W-\bar{W})=W-t(W-E)-\theta E\\
  &=(1-t)W-(\theta-t)E\\
  \end{aligned}
  \]
\end{itemize}
\end{frame}

\begin{frame}
  \frametitle{Optimizing on how much to report}
  Household problem
   \begin{itemize}\setlength\itemsep{3pt}
\item We take the expected utility framework
\[
V= (1-p)u(Y_N) + p u(Y_A) = (1-p)u[(1-t)W+tE] +p u[(1-t)W-(\theta-t)E]
\]
\item Solve the FOC with respect to $E$
\[
\begin{aligned}
  [E]: &\ \  (1-p)u'(Y_N)t - p u'(Y_A)(\theta-t)=0 \\
  & \iff \frac{u'(Y_N)}{u'(Y_A)}= \frac{p}{1-p}\frac{\theta-t}{t}\\
\end{aligned}
\]
\item MRS of evaded income = MC of tax evasion (both relative to audited income)
\end{itemize}
\end{frame}

\begin{frame}
  \frametitle{Comparative statics}
  Comparative statics on $p$ and $\theta$
   \begin{itemize}\setlength\itemsep{3pt}
\item Rise in $p$ means audits are more likely to happen and $\frac{\partial E}{\partial p}<0$
\begin{itemize}\setlength\itemsep{3pt}
  \item Rise in $p$ raises $\frac{p}{1-p}  \Rightarrow\frac{u'(Y_N)}{u'(Y_A)}$ should rise to preserve equality (or $Y_N$ should fall)
  \item Since $Y_N=(1-t)W+tE$ and $Y_A=(1-t)W - (\theta-t)E$ this implies $E\downarrow$
\end{itemize}
\item Rise in $\theta$ increases penalty once detected and $\frac{\partial E}{\partial \theta}<0$
\begin{itemize}\setlength\itemsep{3pt}
  \item Rise in $\theta$ raises $\frac{\theta-t}{t}  \Rightarrow\frac{u'(Y_N)}{u'(Y_A)}$ should rise to preserve equality (or $Y_N$ should fall)
  \item With similar logic as above, $E\downarrow$
\end{itemize}
\end{itemize}
\end{frame}



\begin{frame}
  \frametitle{Comparative statics}
Comparative static on tax rate $t$ is a bit more complicated
\begin{itemize}\setlength\itemsep{3pt}
  \item Substitution effect: Higher tax increases gains from hiding income (without affecting penalty anount) and raises $E$
  \item Income effect: Higher tax makes taxpayer poorer anyway, and if poor are risk averse (of being caught), then they decrease $E$ 
\item This is because penalty $\theta E$ does not scale with $t$
\begin{itemize}
  \item Rising $t$  increases gains from evasion w/o proportionally increasing the cost of detection
\end{itemize}
\end{itemize}\medskip
Adjust the model: Pay taxes on evasion itself $\theta t E$
\begin{itemize}\setlength\itemsep{3pt}
  \item Then $Y_A=(1-t)W + tE(1-\theta)$ and  FOC on tax evasion becomes
  \[
  \frac{u'(Y_N)}{u'(Y_A)}= \frac{p(\theta-1)}{1-p}
  \]
  \item Here, rise in $t$ raises both gains from evasion \textit{and} penalty if caught
 \item Thus, there is no substitution effect, and $E$ unambiguously decreases
\end{itemize}
\end{frame}


\begin{frame}
\frametitle{Higher taxes $\to$ Less evasion?}
Intuitively, higher taxes increase the burden on taxpayers
\begin{itemize}\setlength\itemsep{3pt}
\item More to gain from hiding income
\item Natural expectation: higher $t$ $\Rightarrow$ more evasion
\end{itemize}\medskip
 Previous results shows the opposite: higher $t$ $\Rightarrow$ less evasion
\begin{itemize}\setlength\itemsep{3pt}
\item When penalty is on evaded taxes $\theta t E$, gains and penalties scale equally with $t$
\item $t$ cancels out of the evasion decision entirely
\item Higher $t$ only makes taxpayer poorer, reducing appetite for risky evasion
\end{itemize}\medskip
Empirically, most evidence goes against this theory
\begin{itemize}\setlength\itemsep{3pt}
\item Higher tax rates tend to be associated with more evasion in practice
\item Suggests substitution effect dominates in real settings
\item Or that penalty structures do not follow the this assumption in practice
\end{itemize}
\end{frame}

\begin{frame}
\frametitle{What this means in the context of developing countries}
Developing countries generally have low detection capacities
\begin{itemize}\setlength\itemsep{3pt}
\item In terms of our model, $p$ is low
\item Implying higher $E$ and less tax revenues
\item Example: Informal economies (undetected from formal state)
\end{itemize}\medskip
 They also may have lower $\theta$
\begin{itemize}\setlength\itemsep{3pt}
\item Lack of legal enforcement on deviant behaviors
\item Thus, more evasion (and less tax revenues)
\item Example: Less credible threat of enforcing penalties (failure to monopolize tax)
\end{itemize}\medskip
\end{frame}


\section{Informal economies and taxation}
\subsection{Gordon, Li (2009): Tax structures in developing countries: Many puzzles and a possible explanation}
\begin{transitionframe}
\begin{center}
  \begin{huge}
    \textcolor{NTHU1}{%
  Gordon, Li (2009):\\ Tax structures in developing countries: Many puzzles and a possible explanation
    }
  \end{huge}
\end{center}
\end{transitionframe}

\begin{frame}
\frametitle{Q: Do presence of informal sectors affect choice of tax policies?} 
Major difference in tax structures across developing and developed countries
\begin{itemize}\setlength\itemsep{3pt}
\item Some could be attributed to social preferences over public goods and redistribution
\item But many differences with predicted optimal taxation in generic settings
\begin{itemize}\setlength\itemsep{3pt}
\item These taxation preserves production efficiency under plausible assumptions
\item Rule out tariffs, capital income taxes, differential consumption tax, inflation tax
\end{itemize}
\item Drastic diffrences in developing countries
\begin{itemize}\setlength\itemsep{3pt}
\item Less personal income tax, more tariffs, corporate income tax, and seignorage
\end{itemize}
\end{itemize}\medskip
This paper proposes that existence of informal sector explains different tax structures
\begin{itemize}\setlength\itemsep{3pt}
\item Incorporates enforcement problems in developing countries due to informal economy
\begin{itemize}\setlength\itemsep{3pt}
\item Informal economies: Use cash over financial intermediaries
  \item Taxes on formal economy shifts activities to informal economy $\to$ Optimal tax rate $\Downarrow$
  \item Justifies inflation tax, differential tax across industries, and tariffs
\end{itemize}

\end{itemize}\medskip
\end{frame}

\begin{frame}
\frametitle{Summary of the difference in tax choices across country income} 
\begin{figure}
  \includegraphics[keepaspectratio, width=\textwidth]{gl2009_t1.png}
\end{figure}
\begin{itemize}\setlength\itemsep{3pt}
\item  High corporate income tax, tariffs, seignorage in developing countries
\item Less personal income tax
\item Statutory tax rates are similar across different countries!
\end{itemize}
\end{frame}

\begin{frame}
\frametitle{Tax policy with limited information: Differential excise taxes} 
Characterizing informal economy into a model
\begin{itemize}\setlength\itemsep{3pt}
\item Informal economy not observed by government
\begin{itemize}\setlength\itemsep{3pt}
\item Thus, taxes cannot be levied to such firms
\item So profits for such firms are revenues - wage cost - capital costs
\end{itemize}
\item Formal firms use financial systems and have higher outputs, but subject to tax
\item Firms become formal if net profits exceed those from informal economy
\begin{itemize}\setlength\itemsep{3pt}
\item High excise taxes reduce the net profits $\to$ tax revenue drops
\end{itemize}
\end{itemize}\medskip
Variation across industries
\begin{itemize}\setlength\itemsep{3pt}
\item Sectors with low increase in output from formalizing $\to$ lower tax rates than desired
\item This results in distortion in domestic production, which are offset by tariffs on imports
\item Inflation tax and red tape on lightly taxed sectors may be attractive
\end{itemize}
\end{frame}




\begin{frame}
\frametitle{How to increase benefits of using financial sector} 
Determinants of benefits to using financial sectors
\begin{itemize}\setlength\itemsep{3pt}
\item Capital intensive firms: They need lending to finance purchase of capital
\item If firms operate at wide distance, financial sector reduces the cost
\item High revenue per employee: More cash transaction exposes to internal theft
\begin{itemize}
\item Service sectors tend to be less capital intensive and easily transfer to informal economy
\end{itemize}
\item Firms that do not formalize avoid corporate and excise tax, but pay inflation tax
\item Thus, taxes are kept low for non capital-intensive sectors
\begin{itemize}\setlength\itemsep{3pt}
\item Capital-intensive ones more likely to stick to formal even without intervention
\item This is meant to shift informal economies to formal ones
\end{itemize}

\item This explains high capital income tax in developing countries (vs. income tax)
\item Inflation taxes also raise cost of keeping cash, inducing firms to formalize
\end{itemize}\medskip

\end{frame}

\begin{frame}
\frametitle{Motivating uses of other anti-competitive measures}
Tariffs  
\begin{itemize}\setlength\itemsep{3pt}
\item Higher tariff protection for capital intensive sectors
\begin{itemize}\setlength\itemsep{3pt}
\item Meant to offset distortions due to uneven domestic tax enforcement
\item Specifically, aim is to induce production to a sector more likely to formalize
\end{itemize}
\end{itemize}\medskip
Other anti-competitive measures
\begin{itemize}\setlength\itemsep{3pt}
\item Provide monopoly power to capital-intensive firms
\item Impose tighter regulation or monitoring on informal firms
\begin{itemize}\setlength\itemsep{3pt}
\item Even if revenue from such firms cannot be observed, other activities could be
\item Higher red tape measure on those firms as monitoring device
\end{itemize}
\end{itemize}\medskip
\end{frame}

\begin{frame}
\frametitle{Key takeaways}
Existance of informal firms twist the choice of optimal taxation 
\begin{itemize}\setlength\itemsep{3pt}
\item More capital taxation, tariffs, and inflation taxes
\item This is determined by the returns to using financial sectors
\begin{itemize}\setlength\itemsep{3pt}
\item If financial sectors provide benefits, less incentives to join informal economy
\item As developing countries have less developed finance sectors, high informality
\end{itemize}
\item Thus, models should incorporate possiblity of informal economies
\end{itemize}\medskip
However, leaving informal economies as is is undesirable
\begin{itemize}\setlength\itemsep{3pt}
\item Weakens capacity to provide essential public goods
\item Resource misallocation$\Uparrow$ and innovation$\Downarrow$ (cost of hiding leave little for innovation)
\item Thus, efforts to induce firms to formalize through third-party information, legal measures, and expansion of electronic filing (Benhassine et al 2018)
\end{itemize}
\end{frame}

\subsection{Benhassine, Mckenzie, Pouliquen, Santini (2018): Does Inducing Informal Firms to Formalize Make Sense? Evidence from Benin }
\begin{transitionframe}
\begin{center}
  \begin{huge}
    \textcolor{NTHU1}{%
Benhassine, Mckenzie, Pouliquen, Santini (2018): \\Does Inducing Informal Firms to Formalize Make Sense? Evidence from Benin
    }
  \end{huge}
\end{center}
\end{transitionframe}


\begin{frame}
\frametitle{Q: Can governments induce informal firms to formalize?}
Many small firms in the developing country are informal
\begin{itemize}\setlength\itemsep{3pt}
\item In some cases, their estimated size is equivalent to 70\% of GDP
\item These could be very problematic
\begin{itemize}\setlength\itemsep{3pt}
\item Informal firms cannot access formal finance sector, public contracts etc 
\item Lost tax revenues for the government
\end{itemize}
\item Thus, reforms to mandate business entry regulation is in place in many places
\end{itemize}\medskip
This paper uses randomized experiment to see what can be done to induce formalization
\begin{itemize}\setlength\itemsep{3pt}
\item Involves 3,600 firms in Benin
\item Introduction of entreprenant legal status, a simplified formalization regime
\item Estimates benefits for firms and costs of the intervention
\item Also provides suggestion on who is most likely to formalize
\end{itemize}
\end{frame}


\begin{frame}
\frametitle{Institutional Context in Benin}
Entreprenant status reform in Benin
\begin{itemize}\setlength\itemsep{3pt}
\item Came into effect in 2011
\item Aim to facilitate formalization through simplified procedures
\item Costs also reduced by 70-80\%
\item As these were rolled out suddenly, many firms seemed to have been unaware
\end{itemize}\medskip
What happens when firms formalize (vs informal)
\begin{itemize}\setlength\itemsep{3pt}
\item Get tax ID and exempt for taxation for one tax year (and 3-year reduced taxation)
\item Informal firm addresses unknown to govt (so difficult to collect tax)
\end{itemize}
\end{frame}

\begin{frame}
\frametitle{Experiment design}
Treatment across 3595 firms
\begin{itemize}\setlength\itemsep{3pt}
\item A: In-person visits with verbal explanation and assistance in paperworks
\item B: Help open a bank account and provide training in return for formalizing
\item C: Provide tax mediation service and tax preparation service
\item Treatment1 got A, Treatment2 got A+B, Treatment 3 got A+B+C 
\item The experiment data is merged with admin data on formalization and firm surveys
\end{itemize}\medskip
The regression follows an OLS format, given randomized assignment
\[
y_i=\alpha + \beta_1 T1_i+ \beta_2 T2_i ++\beta_3 T3_i +X_i +e_i
\]\vspace{-20pt}
\begin{itemize}\setlength\itemsep{3pt}
\item Follow-up surveys: Form a panel of firm-years
\[
\begin{aligned}
y_{it}&=\alpha + b_1 (T1_i\cdot F1)+ b_2 (T2_i\cdot F1)+b_3 (T3_i\cdot F1)\\
&+ c_1 (T1_i\cdot F2)+ c_2 (T2_i\cdot F2)+c_3 (T3_i\cdot F2)+\pi Y_{i0} + \gamma M_{i0} +X_i\delta +e_i
\end{aligned}
\]
\end{itemize}
\end{frame}


\begin{frame}
\frametitle{More firms to formalize}
\begin{figure}
\centering
\includegraphics[keepaspectratio, width=0.8\textwidth]{bmps2019_f1.png}
\end{figure}
\end{frame}


\begin{frame}
\frametitle{But limited effect on take-up of further programs}
\begin{figure}
\centering
\includegraphics[keepaspectratio, width=0.9\textwidth]{bmps2019_t1.png}
\end{figure}
\end{frame}

\begin{frame}
\frametitle{And limited impact on firm performance}
\begin{figure}
\centering
\includegraphics[keepaspectratio, width=0.9\textwidth]{bmps2019_t2.png}
\end{figure}
\end{frame}

\begin{frame}
\frametitle{Cost-benefit analysis and takeaways}
Cost-benefit analysis
\begin{itemize}\setlength\itemsep{3pt}
\item Total costs in the 2 years: Equivalent to US\$ 620,000
\item Costs per formalization: Range from \$1,237 -- \$ 2,217
\item This is more than 10 times of the profits from formalization
\item Recouping costs through taxes: Takes estimated 20 years
\item Better targeting could help: More educated firmowners more likely to formalize
\end{itemize}\medskip
Takeaways
\begin{itemize}\setlength\itemsep{3pt}
\item Formalization can bring in additional tax revenues in theory
\item But implementation is costly: Takes time to recoupe costs of additinal interventions
\item What now? Target those who are more likely to formalize?
\end{itemize}
\end{frame}


\section{Conflicts and implementation of tax authority}
\subsection{Sanchez de la Sierra (2020): On the Origins of the State: Stationary Bandits and Taxation in Eastern Congo }
\begin{transitionframe}
\begin{center}
  \begin{huge}
    \textcolor{NTHU1}{%
Sanchez de la Sierra (2020):\\ On the Origins of the State: Stationary Bandits and Taxation in Eastern Congo
    }
  \end{huge}
\end{center}
\end{transitionframe}

\begin{frame}
\frametitle{Q: How do stationary bandits mimick essential functions of the state?}
Monopoly of violence, taxation, and law enforcement is an essential function of state
\begin{itemize}\setlength\itemsep{3pt}
\item But even in developed countries, these are modern concepts
\item In what environment do these functions emerge?
\item However, many quantitative data exist \textit{since} the states are formed
\item so what explains \textit{the birth} of states?
\end{itemize}\medskip
This paper investigates this question using a natural experiment
\begin{itemize}\setlength\itemsep{3pt}
\item Uses Democratic Republic of Congo (DRC) as a context
\begin{itemize}\setlength\itemsep{3pt}
\item Vacuum of the formal state presence (DRC lost de facto control of part of its territory)
\item Leaving armed actors to monopolize violence, taxation, and other essential services
\end{itemize}
\item Motives of armed actors determined by external changes in resource prices
\item Rise of state function related to which types of resource shock is in place
\end{itemize}
\end{frame}


\begin{frame}
\frametitle{Conceptualizing state functions}
When are states likely to emerge?
\begin{itemize}\setlength\itemsep{3pt}
\item States monopolize violence to repress threats to the property of the governed
\item They then monopolize taxation to finance their operations
\item The potential tax revenues motivate armed actors to protect welfare of the population
\end{itemize}\medskip
Motivating hypotheses
\begin{itemize}\setlength\itemsep{3pt}
\item When do state emerge and develop state capacity?
\begin{itemize}
\item Likely when an output that can be easily taxed exists
\item The output themselves should promise significant returns for the armed actors
\end{itemize}
\item What is the net benefit for the population?
\begin{itemize}
\item If investments spurred by protection $>$ disincentives created by taxing production
\item If taxing output (suppliers) is difficult, states may choose to target consumption
\end{itemize}
\item In production economies: State emerges with high output prices at production sites
\item In consumption economies: State emerges where the outputs are being consumed
\end{itemize}\medskip
\end{frame}


\begin{frame}
\frametitle{Setting and data in DRC}
Setting: Eastern DRC in early 2000s
\begin{itemize}\setlength\itemsep{3pt}
\item Mineral extraction in eastern DRC (lack of formal state) promises high returns
\item Many militias (regional, external actors, DRC army) have incentives to take over
\item Demand shock for different minierals
\begin{itemize}\setlength\itemsep{3pt}
\item Coltan: Heavy, demand shock driven by announcement of PS2
\item Gold: Light/concealable, safe haven of values during recession
\end{itemize}
\end{itemize}\medskip
Data and analysis
\begin{itemize}\setlength\itemsep{3pt}
\item Mining site and local history (outputs) compiled by surveyors on the ground
\item Uses the following municipality-year level regression
\[
Y_{jt}=\beta_t +\alpha_j + \gamma_{\text{coltan}}C_jp_{\text{coltan},t}^{\text{US}}+ \gamma_{\text{gold}}G_jp_{\text{gold},t}^{\text{US}} + e_{jt}
\]
\end{itemize}\medskip
\end{frame}


\begin{frame}
  \frametitle{State emerges differently in production and consumption economies}
\begin{figure}
  \centering
  \includegraphics[keepaspectratio, width=1\textwidth]{s2020_t1.png}
\end{figure}
\end{frame}


\begin{frame}
  \frametitle{Coltan: Benefits locations without stationary bandits}
\begin{figure}
  \centering
  \includegraphics[keepaspectratio, width=1\textwidth]{s2020_t2.png}
\end{figure}
\end{frame}


\begin{frame}
  \frametitle{Gold: Benefits locations with stationary bandits}
\begin{figure}
  \centering
  \includegraphics[keepaspectratio, width=1\textwidth]{s2020_t3.png}
\end{figure}
\end{frame}

\begin{frame}
\frametitle{Takeaways and discussion}
Summaries of key findings
\begin{itemize}\setlength\itemsep{3pt}
\item Armed groups sometimes form monopoly of violence, taxation, and protection
\item Especially if outputs with high returns that can be expropriated are present
\item In turn, they form essential state functions that sometimes benefitting the well-being
\end{itemize}\medskip
Remaining questions
\begin{itemize}\setlength\itemsep{3pt}
\item Sometimes the population would rather support armed groups over the government
\begin{itemize}
  \item e.g. Mexico, Colombia, or any places where organized criminal groups provide services
\end{itemize} 
\item But that is not always the case
\item So how do armed groups legitimize their presence?
\end{itemize}\medskip
\end{frame}


\subsection{Henn, Mugaruka, Ortiz, Sánchez de la Sierra, Wu (2025): Monopoly of Taxation without a Monopoly of Violence: The Weak State's Trade-off from Taxation}
\begin{transitionframe}
\begin{center}
  \begin{huge}
    \textcolor{NTHU1}{%
  Henn, Mugaruka, Ortiz, Sánchez de la Sierra, Wu (2025):\\ Monopoly of Taxation without a Monopoly of Violence: The Weak State's Trade-off from Taxation
    }
  \end{huge}
\end{center}
\end{transitionframe}

\begin{frame}
\frametitle{Q: Tradeoffs between monopoly of taxation \& monopoly of violence?} 
Weak states often fail to exert de facto control over their jurisdictions
\begin{itemize}\setlength\itemsep{3pt}
\item Non-state actors (rebels) take over, who sometimes provide stability over 
\item If these areas are economically valuable, state loses lot of fiscal revenue
\item So is it beneficial for state to take formal authority over those lands?
\begin{itemize}\setlength\itemsep{3pt}
\item Regain lost fiscal revenues
\item But could trigger inadvertent response from rebels (e.g. plundering)
\end{itemize}
\end{itemize}\medskip
This paper investigates the effect of state-building attempts to increase tax revenues
\begin{itemize}\setlength\itemsep{3pt}
\item Military powers to reclaim land and replace rebels
\item Negotiate with rebels (often with concession to the rebels)
\item Leverage subnational evidence in Democratic Republic of Congo (DRC) 
\end{itemize}
\end{frame}

\begin{frame}
\frametitle{State-building efforts to drive rebels out of DRC} 
State absence in Eastern DRC (and postcolonial Africa in general)
\begin{itemize}\setlength\itemsep{3pt}
\item Many parts of DRC (East) were under armed group control since 1998
\item Most notable of these organizations is FDLR (Forces de Libération du Rwanda)
\end{itemize}\medskip
Military operation and peace treaty with rebels
\begin{itemize}\setlength\itemsep{3pt}
\item Largest peace treaty: Sun City peace treaty in 2004
\begin{itemize}\setlength\itemsep{3pt}
\item Various armed groups integrated into Congolese army
\item Allowed government to regain half the country without triggering violence
\item Other peace agreements had smaller goals and lower credibility
\end{itemize}
\item Largest military operation: Kimia II (March 2009)
\begin{itemize}\setlength\itemsep{3pt}
\item Joint operation with UN, national army, and Rwanda
\item Targeted FDLR presiding and taxing over Basile Chiefdom region since 2005
\item Although defeated, FDLR launched theft operations from the nearby hilly forest region
\end{itemize}
\end{itemize}
\end{frame}

\begin{frame}
\frametitle{Data and empirical strategy} 
Constructing the data
\begin{itemize}\setlength\itemsep{3pt}
\item Armed group presence: village-level data through interviews with former combatants
\begin{itemize}
\item Perpetrators' affiliation, motivation (pilage/retaliate/etc), actions taken
\end{itemize}
\end{itemize}\medskip
Effects of Kimia II and Sun City Peace Treaty
\begin{itemize}\setlength\itemsep{3pt}
\item Kimia II: FDLR vs non-FDLR villages pre (`05-`09) and post (`10-`12) in Basile Chiefdom
\item Peace treaty: Integrated villages vs others in the region before and after `02
\[
Y_{i,t} = \alpha_i + \beta_t + \sum_{k=-4}^{k=3}\beta_k \text{Treat}_i \times 1[t=t_0+k] + e_{i,t}
\]\vspace{-20pt}
\begin{itemize}
  \item $t_0=2009$ for Kimia II effects, $2002$ for peace treaty effects
\end{itemize}
\end{itemize}
\end{frame}

\begin{frame}
\frametitle{Kimia II: FDLR increased violence afterwards} 
\begin{figure}
\includegraphics[keepaspectratio,width=0.9\textwidth]{hmosw2025_f1.png}  
\end{figure}
\end{frame}

\begin{frame}
\frametitle{Kimia II: Driven by increase in pillages by FDLR} 
\begin{figure}
\includegraphics[keepaspectratio,width=0.9\textwidth]{hmosw2025_t1.png}  
\end{figure}
\end{frame}

\begin{frame}
\frametitle{Peace treaty: Rise of new armed groups since treaty} 
\begin{figure}
\includegraphics[keepaspectratio,width=0.7\textwidth]{hmosw2025_f2.png}  
\end{figure}
\end{frame}

\begin{frame}
\frametitle{Takeaways and discussions} 
What (doesn't) explain the effects
\begin{itemize}\setlength\itemsep{3pt}
\item Actual violence increased (as opposed to \textit{reported} violence)
\item Violence by other actors do not increase following Kimia II
\begin{itemize}
\item Thus when incentive to protect villages (for tax revenues) vanish, violence increase
\end{itemize}
\item Treaty raised revenues, but limited commitment ability led to questioning of legitimacy
\end{itemize}\medskip
Takeaways and questions
\begin{itemize}\setlength\itemsep{3pt}
\item Kimia II: Monopolize violence (rebels out), but lost taxation due to ensuing violence
\item Peace treaty: Expand tax revenues but no longer monopolize violence
\item If faced with these trade-offs what should the states do?
\begin{itemize}\setlength\itemsep{3pt}
  \item Asserting force desirable if state has more capacity to prevent crime 
  \item Bargaining potentially desirable if states can commit to not negotiate with future rebels and have resource to integrate rebels
\end{itemize}
\end{itemize}
\end{frame}

 
\begin{frame}
\frametitle{Unanswered questions} 
Informal economies reduce potential for tax revenues
\begin{itemize}\setlength\itemsep{3pt}
\item Thus, states have less information about tax payers
  \item This induces developing countries to choose different types of taxes
\begin{itemize}\setlength\itemsep{3pt}
\item Those that are easily detectable and difficult to avoid
\item Trade taxes, corporate income tax, consumption tax 
\end{itemize}
\item But formalizing them is also costly
\item Then, how to increase amount of information (e.g. new technologies?)
\end{itemize}\medskip
Failure to monopolize taxes
\begin{itemize}\setlength\itemsep{3pt}
\item States sometimes compete with non-state actors (armed groups / local chiefs)
  \item Tax administration is limited in scope
\item Any way to incorporate local knowledge?
\end{itemize}
\end{frame}
%%%%%%%%%%%
\end{document}
